% =============================================================================
% INCOHERENTIA XII: DE SCRIPTURIS ET FORMIS LITTERARUM IN ARCHIVO FONTIUM
% Status: DRAFTED
% Typographic state: BYOBU DRIFT BEGINS HERE. Arabic page counter runs ahead.
% This chapter is the scholastic font excursus — rendered primarily in Amiri
% (classical authority), but with SGA and Cursive SGA intrusions that undermine
% the authority being claimed. The reader is given a theory of the system at
% the exact moment the system ceases to obey it.
% Pagestyle: byobu-drift (set in main.tex before this chapter)
% =============================================================================

\chapter*{Incoherentia XII: De Scripturis et Formis Litterarum in Archivo Fontium}
\addcontentsline{toc}{chapter}{XII. De Scripturis et Formis Litterarum}

\setstretch{1.0}

{\AmiriFont
In the middle years of the compilation, when the manuscript had already begun
to exhibit signs of internal divergence, the Editor records a discovery not
of narrative content, but of the instruments through which narrative itself
is made manifest. These were not found in the body of the text, nor in its
marginal apparatus, but in a subdirectory of the archive, designated simply
as \texttt{fonts/}, as though the naming itself sought to conceal the gravity
of its contents beneath an ordinary sign.

\medskip

Within this directory were gathered a multiplicity of scriptural forms, each
encoded in files bearing the extensions \texttt{.ttf} and others more obscure,
suggesting not merely variation of appearance but divergence of ontological
regime. The Editor, proceeding in the manner of the scholastics, did not at
first distinguish between them by aesthetic quality, but by function,
comportment, and the manner in which each seemed to impose constraints upon
the text it rendered.

\medskip

Foremost among these was the family designated \texttt{Sga-Regular}, present
in several instantiations. This script, which the Editor identifies with the
Standard Galactic hand, exhibits a character of rigid clarity and geometric
finality. Its glyphs appear less written than \sga{inscribed}, as though each
letter were the projection of a more primitive symbolic order. When passages
are rendered in this form, ambiguity diminishes, yet at the same time the text
becomes resistant to paraphrase, as if each sign were bound to a singular and
irreducible \sga{meaning}. It is therefore associated with the underlying
structure of the manuscript, that which persists even when interpretation
fails.

\medskip

In contrast, the script preserved as \texttt{CursiveGalactic-Regular.ttf} is
of a markedly different nature. Its lines flow and interweave, and its
characters seem less discrete than continuous, as though each were the residue
of motion rather than the product of intention. \cga{Passages rendered in this
hand exhibit a tendency toward instability, with meanings that drift and
recombine.} The Editor therefore assigns to this script the domain of
transition, wherein the text has not yet settled into fixed form, and where
interpretation must proceed by following trajectories rather than identifying
points.

\medskip

A third principal form is encountered in \texttt{NovaMonoStandardGalactic.ttf},
which unites the monospaced discipline of Nova Mono with the glyphic repertoire
of the Galactic script. Here, each character occupies an equal measure of space,
producing a lattice-like regularity that suggests computation or enumeration.
Yet within this grid, the foreign shapes of the Galactic letters introduce a
tension between uniform spacing and non-uniform meaning. The Editor remarks that
texts rendered in this form often present themselves as stable, yet conceal
within their regularity a deeper incoherence, detectable only when one attends
to the mismatch between alignment and sense.

\medskip

Among the ancillary forms, the Amiri family presents a script of classical
dignity, reminiscent of juridical or theological treatises. Its curves and
proportions evoke continuity with earlier traditions, and passages rendered
therein acquire an air of authority, whether or not such authority is warranted.
It is therefore frequently employed in sections where the manuscript feigns
coherence, adopting the voice of established doctrine even as it quietly
undermines it.

\medskip

The scripts designated \texttt{Cheiro-Regular} and \texttt{Clypto-Regular}
appear to belong to a more enigmatic lineage. Their forms suggest systems of
encoding or cipher, and indeed the Editor notes that passages rendered in these
hands often resist immediate comprehension, yielding their sense only through
repeated inspection or cross-reference. They are thus associated with the hidden
strata of the text, where meaning is not absent but \cga{deferred}.

\medskip

Particular attention is given to \texttt{Logico\_philosophicus-Regular.ttf},
whose very name invokes the ambition of total systematization. The glyphs here
are austere, almost diagrammatic, and passages composed in this script tend
toward propositions, enumerations, and structures that aspire to closure. Yet,
as the Editor observes, such closure is never complete; \sga{the very rigidity
of the form reveals the limits of the system it seeks to impose}.

\medskip

Likewise, \texttt{Systada-Regular.ttf}, described in the accompanying log as
terminal-friendly, introduces a script adapted to mechanical environments. Its
characters are optimized for legibility under constraint, and texts rendered
therein often take on the aspect of commands, logs, or outputs. In these
passages, the manuscript appears to address not a reader but a process, as
though the act of reading had been displaced into execution.

\medskip

Finally, there are forms such as \texttt{shapeform.ttf} and \texttt{dactyl.ttf},
whose glyphs exhibit experimental or incomplete qualities. These are treated by
the Editor not as finished scripts but as traces of a generative process,
remnants of an attempt to discover new alphabets rather than merely deploy
existing ones. Their presence suggests that the manuscript is not only written
but continually rewriting the very means of its own \cga{inscription}.

\medskip

From these observations, the Editor concludes that the fonts are not
interchangeable vessels for a preexisting text, but active participants in its
formation. Each imposes constraints, affords possibilities, and alters the space
of meanings that can be realized. To read the manuscript without regard to the
script in which it is rendered is therefore to ignore a primary dimension of its
\sga{structure}.

\medskip

And it is further noted, though not proven, that in certain passages the
transformation from one script to another does not merely represent a change in
presentation, but corresponds to a transition in the state of the narrative
itself, as though the text, encountering a failure of coherence, were compelled
to seek refuge in another form of \cga{inscription}.

} % end AmiriFont

\medskip

\begin{AR}
وقد لاحظ المحرر أن الخط ليس وعاءً للمعنى، بل هو المعنى نفسه في صورته الأولى.
\end{AR}
